\section{Результаты}

\subsection*{Количество документов}

Объём корпуса оценён по счётчикам алфавитных категорий русской Википедии
(таблица~\ref{tab:cats}). Вложенные подкатегории не обходились, поэтому это нижняя оценка.

\begin{table}[h]
\centering
\begin{tabular}{|l|r|}
\hline
\textbf{Категория} & \textbf{Статей} \\
\hline
Фильмы по алфавиту & 38\,118 \\
Актёры по алфавиту & 15\,563 \\
Сценаристы по алфавиту & 6\,735 \\
Кинорежиссёры по алфавиту & 5\,502 \\
Телесериалы по алфавиту & 5\,148 \\
Мультфильмы по алфавиту & 4\,322 \\
\hline
\textbf{Итого} & \textbf{75\,388} \\
\hline
\end{tabular}
\caption{Объём кинотематики в русской Википедии}
\label{tab:cats}
\end{table}

Требование в 30\,000 документов перекрыто больше чем вдвое одним только первым
источником.

У film.ru статьи собраны в три раздела-списка по 28 статей на странице: рецензии ---
324 страницы, подборки --- 187, интервью --- 34. Номер последней страницы найден
перебором: следующая за ней отвечает 404. Получается около 15 тысяч статей. Вместе с
Википедией корпус заведомо больше требуемых 30 тысяч, даже с учётом того, что часть
статей окажется заглушками.

\subsection*{Размеры документов}

Измерения по 13 сохранённым страницам приведены в таблице~\ref{tab:sizes}. Под
\enquote{сырым} размером понимается HTML-файл целиком, под текстом --- то, что осталось
после выделения. Доля --- объём текста в байтах к размеру страницы; оба источника
отдают UTF-8, так что байты считаются одинаково.

\begin{table}[h]
\centering
\begin{tabular}{|l|r|r|r|r|}
\hline
\textbf{Документ} & \textbf{Сырой, КБ} & \textbf{Символов} & \textbf{Слов} & \textbf{Доля} \\
\hline
\multicolumn{5}{|l|}{\textit{ru.wikipedia.org}} \\
\hline
Во все тяжкие (сериал)      & 1152,3 & 161\,473 & 18\,159 & 18,9\,\% \\
Янковский, Олег Иванович    & 1011,4 & 135\,005 & 17\,758 & 21,5\,\% \\
Тарковский, Андрей Арс.     & 706,1 & 65\,397 & 9\,568 & 15,3\,\% \\
Криминальное чтиво          & 633,4 & 42\,958 & 6\,866 & 10,3\,\% \\
Сталкер (фильм)             & 523,3 & 45\,539 & 6\,967 & 14,6\,\% \\
13 дней во Франции          & 119,9 & 745 & 133 & 1,0\,\% \\
24 недели                   & 100,4 & 873 & 120 & 1,3\,\% \\
\hline
\multicolumn{5}{|l|}{\textit{www.film.ru}} \\
\hline
Интервью о кинопрокате               & 122,1 & 16\,115 & 2\,456 & 23,4\,\% \\
Рейтинг сериалов М. Флэнагана        & 137,7 & 12\,219 & 1\,750 & 15,7\,\% \\
Обзор сериала \enquote{Полураспад}   & 116,8 & 7\,206 & 980 & 11,0\,\% \\
Рецензия \enquote{Последнее дело Холмса} & 107,9 & 5\,182 & 741 & 8,6\,\% \\
Рецензия на фильм \enquote{Животные} & 111,7 & 4\,579 & 628 & 7,3\,\% \\
\textit{Карточка фильма}             & 149,1 &     369 &     52 &  0,4\,\% \\
\hline
\end{tabular}
\caption{Размеры \enquote{сырых} документов и выделенного текста}
\label{tab:sizes}
\end{table}

Средние значения сведены в таблицу~\ref{tab:avg}. Карточка фильма film.ru в средние
не входит: документом она не считается и приведена только для сравнения.

\begin{table}[h]
\centering
\begin{tabular}{|l|r|r|}
\hline
\textbf{Показатель} & \textbf{Википедия} & \textbf{film.ru} \\
\hline
Документов в выборке                          &       7 &       5 \\
Средний размер \enquote{сырого} документа, КБ &   606,7 &   119,2 \\
Средний объём текста, символов                & 64\,570 &  9\,060 \\
Средний объём текста, слов                    &  8\,510 &  1\,311 \\
Средняя доля полезного текста                 & 11,9\,\% & 13,2\,\% \\
\hline
\end{tabular}
\caption{Средние характеристики документа}
\label{tab:avg}
\end{table}

Самое заметное здесь --- доля полезного текста: у Википедии в среднем 12\,\%, а у
статей-заглушек всего 1\,\%. Причина в оформлении: тема встраивает в каждую страницу
около полумегабайта меню и скриптов независимо от длины статьи. Поэтому судить об объёме
будущего индекса по весу скачанных данных нельзя.

У film.ru картина другая: страницы в пять раз легче, а статьи короче, но ровнее ---
от 600 до 2\,400 слов, заглушек в выборке нет. Доля текста при этом даже выше, чем у Википедии.
Карточка фильма весит больше любой статьи, а текста в ней 52 слова --- поэтому
карточки и не взяты в корпус.

\pagebreak
